% XeLaTeX can use any Mac OS X font. See the setromanfont command below.
% Input to XeLaTeX is full Unicode, so Unicode characters can be typed directly into the source.

% The next lines tell TeXShop to typeset with xelatex, and to open and save the source with Unicode encoding.

%!TEX TS-program = xelatex
%!TEX encoding = UTF-8 Unicode

\documentclass[12pt,twoside]{article}

\usepackage{geometry}                % See geometry.pdf to learn the layout options. There are lots.
\geometry{a4paper}                   % ... or a4paper or a5paper or ... 
%\geometry{landscape}                % Activate for for rotated page geometry
%\usepackage[parfill]{parskip}    % Activate to begin paragraphs with an empty line rather than an indent
\usepackage{amssymb}
\usepackage{todonotes}
\usepackage{booktabs}


% Will Robertson's fontspec.sty can be used to simplify font choices.
% To experiment, open /Applications/Font Book to examine the fonts provided on Mac OS X,
% and change "Hoefler Text" to any of these choices.

\usepackage{fontspec,xltxtra,xunicode}
\defaultfontfeatures{Mapping=tex-text}
\setromanfont[Mapping=tex-text]{Hoefler Text}
\setsansfont[Scale=MatchLowercase,Mapping=tex-text]{Optima}
\setmonofont[Scale=MatchLowercase]{Andale Mono}

%%%% Graphics
\usepackage{graphicx}
\usepackage{tikz}
\definecolor{unilublue}{RGB}{55,149,218}
\definecolor{sntred}{RGB}{219,46,27}
\definecolor{sntpurple}{RGB}{86,30,130}
%\definecolor{sntblue}{RGB}{50,130,207}
\definecolor{sntblue}{RGB}{55,149,218}

\pgfdeclareimage[width=30mm]{logo-snt}{logos/logo-snt}
\pgfdeclareimage[width=30mm]{logo-uni-lu}{logos/logo-uni-lu}

%%%% Fancy
\usepackage{fancyhdr}

%%%% Increase page length
\addtolength{\textheight}{1in}

%%%% Sections
\usepackage{sectsty}
\allsectionsfont{\sffamily}

%%%% Title
\newcommand{\titleone}{\textsf{Mutation Testing Survey:}}
\newcommand{\titletwo}{\textsf{FAQAS}}

\newcommand{\todoinline}[1]{\todo[color=orange,inline]{ \textbf{Note}: #1 }}

\begin{document}

\pagestyle{fancy}
\renewcommand{\sectionmark}[1]{\markright{\textit{#1}}}

\renewcommand{\headrulewidth}{2pt}% 2pt header rule
\renewcommand{\headrule}{\hbox to\headwidth{%
  \color{sntblue}\leaders\hrule height \headrulewidth\hfill}}

\fancyhf{}

%\lhead{\fancyplain{}{\setlength{\unitlength}{1mm}
%\begin{picture}(0,0)
%\put(0,-4){\includegraphics[width=50pt]{logos/logo-snt}}
%\end{picture}}} 

\lhead{\fancyplain{}{\textit{Mutation Testing Survey}}}

\rhead{\fancyplain{}{\rightmark }}

\fancyfoot[C]{%
\begin{tikzpicture}[remember picture,overlay]
\path [fill=sntred]    ([xshift=88pt,yshift=20pt]current page.south west) rectangle
                       ([xshift=229pt,yshift=50pt] current page.south west);
\path [fill=sntpurple] ([xshift=229pt,yshift=20pt] current page.south west) rectangle
                       ([xshift=370pt,yshift=50pt] current page.south west);
\path [fill=sntblue]   ([xshift=370pt,yshift=20pt] current page.south west) rectangle
                       ([xshift=510pt,yshift=50pt] current page.south west);
\end{tikzpicture}
}

\fancyfoot[RO]{
\begin{tikzpicture}[remember picture,overlay]
\node [circle, ultra thick, fill=white, draw=sntblue] at ([xshift=530pt,yshift=35pt] current page.south west) {\thepage};
\end{tikzpicture}}

\fancyfoot[LE]{
\begin{tikzpicture}[remember picture,overlay]
\node [circle, ultra thick, fill=white, draw=sntblue] at ([xshift=65pt,yshift=35pt] current page.south west) {\thepage};
\end{tikzpicture}}

\thispagestyle{empty}

\begin{tikzpicture}[remember picture,overlay]
\path [fill=sntred]    ([xshift=30pt,yshift=20pt]current page.south west) rectangle
                       ([xshift=210pt,yshift=50pt] current page.south west);
\path [fill=sntpurple] ([xshift=210pt,yshift=20pt] current page.south west) rectangle
                       ([xshift=390pt,yshift=50pt] current page.south west);
\path [fill=sntblue]   ([xshift=390pt,yshift=20pt] current page.south west) rectangle
                       ([xshift=570pt,yshift=51pt] current page.south west);
\path [fill=unilublue] ([xshift=30pt,yshift= 50pt] current page.south west) --
                       ([xshift=570pt,yshift= 50pt] current page.south west)
                       [rounded corners=20pt] --
                       ([xshift=570pt,yshift=740pt] current page.south west)
                       [sharp corners] --
                       ([xshift=30pt,yshift=740pt] current page.south west);
\node [fill=white,rounded corners=0pt,inner xsep=6pt,inner ysep=3pt]
      at ([xshift=523pt,yshift=120pt] current page.south west)
      {\pgfuseimage{logo-uni-lu}};

\node [fill=white,rounded corners=2pt,inner xsep=6pt,inner ysep=3pt]
      at ([xshift=520pt,yshift=780pt] current page.south west)
      {\pgfuseimage{logo-snt}};

%\node [circle, fill=white, draw=sntblue] at ([xshift=550pt,yshift=35pt] current page.south west) {a};

\node[draw=none,fill=none,right] at (-1, -7){\color{white}\LARGE\bf\titleone};
\node[draw=none,fill=none,right] at (-1, -8){\color{white}\LARGE\bf\titletwo};
\node[draw=none,fill=none,right] at (-1, -9){\color{white}\LARGE\bf\titlethree};

\node[draw=none,fill=none,right] at (-1, -12){\color{white}\Large\textsf{O. Cornejo, F. Pastore, E. Viganò}};
\node[draw=none,fill=none,right] at (-1, -13){\color{white}\Large\textsf{Interdisciplinary Centre for Security, Reliability and Trust}};
\node[draw=none,fill=none,right] at (-1, -14){\color{white}\Large\textsf{University of Luxembourg}};
\node[draw=none,fill=none,right] at (11, -16){\color{white}\textsf{ITT-1-9873-ESA-FAQAS-SUTP}};
\node[draw=none,fill=none,right] at (11, -17){\color{white}\Large\textsf{Issue 4, Rev. 1}};
\node[draw=none,fill=none,right] at (11, -18){\color{white}\Large\textsf{\today}};
\node[draw=none,fill=none,right] at (-1, -24){\color{white}\tiny\textsf{EUROPEAN SPACE AGENCY. CONTRACT REPORT.}};
\node[draw=none,fill=none,right] at (-1, -24.3){\color{white}\tiny\textsf{The work described in this report was done under ESA contract. Responsibility for the contents resides in the author or organisation that prepared it.}};
\node[draw=none,fill=none,right] at (-1, -24.8){\color{white}\tiny\textsf{The copyright in this document is vested in the University of Luxembourg.}};
\node[draw=none,fill=none,right] at (-1, -25.1){\color{white}\tiny\textsf{This document may only be reproduced in whole or in part, or stored in a retrieval system,or transmitted in any form, or by any means electronic,}};
\node[draw=none,fill=none,right] at (-1, -25.4){\color{white}\tiny\textsf{mechanical, photocopying or otherwise, either with the prior permission of the University of Luxembourg or in accordance with the terms of ESTEC Contract No. 4000128969/19/NL/AS.}};

\node[draw=none,fill=none,right] at (-1, -24){\color{white}\tiny\textsf{EUROPEAN SPACE AGENCY. CONTRACT REPORT.}}; 
\node[draw=none,fill=none,right] at (-1, -24.3){\color{white}\tiny\textsf{The work described in this report was done under ESA contract. Responsibility for the contents resides in the author or organisation that prepared it.}};
\node[draw=none,fill=none,right] at (-1, -24.8){\color{white}\tiny\textsf{The copyright in this document is vested in the University of Luxembourg.}};
\node[draw=none,fill=none,right] at (-1, -25.1){\color{white}\tiny\textsf{This document may only be reproduced in whole or in part, or stored in a retrieval system,or transmitted in any form, or by any means electronic,}}; 
\node[draw=none,fill=none,right] at (-1, -25.4){\color{white}\tiny\textsf{mechanical, photocopying or otherwise, either with the prior permission of the University of Luxembourg or in accordance with the terms of ESTEC Contract No. 4000128969/19/NL/AS.}};


\end{tikzpicture}

\newpage


% !TEX root = MutationTestingSurvey.tex

\subsection{Equivalent Mutants}
\label{sec:opt:equivalent}

Equivalent mutants are mutants that behave as the original program, they are semantically equivalent to the original version despite being syntactically different. Equivalent mutants are an actual problem that can lead mutation testing to be infeasible, Schuler and Zeller \cite{schuler2013covering} showed that 45\% of not killed mutants are equivalent, and that manually checking if a mutant is equivalent can take on average up to 15 minutes. Even though, the process of identifying equivalent mutants has been defined as an undecidable problem \cite{madeyski2013overcoming}, the research has cope with heuristics to tackle this issue.

According to the literature \cite{madeyski2013overcoming}, the equivalent mutant heuristics can be classified in two different groups: (1) detecting equivalent mutant techniques, and (2) reducing equivalent mutant techniques.

\subsubsection{Detecting Equivalent Mutant Techniques}

The first group aims to detect and discard the equivalent mutants during the mutation process. In this direction, Papadakis et al. \cite{papadakis2015trivial, kintis2017detecting,papadakis2019mutation} has developed the Trivial Compiler Optimisation technique, the approach relies on the idea that compiler optimisations transforms mutants to the optimised version, thus when the original program can be transformed by an optimisation to one of its mutants, then the mutant is defined equivalent.
Recent studies \cite{papadakis2015trivial} show that the Trivial Compiler Optimisation is able to detect approximately up to 30\% of the equivalent mutants.

A different approach for detecting equivalent mutants is to formulate the problem as a constraint satisfaction problem. The main idea is to analyze the path condition of a mutant, in which the mutant is defined equivalent if and only if the input constraint is unsatisfiable. Offutt et al. \cite{offutt1996detecting,offutt1997automatically} carried out an experimental evaluation on 11 Fortran subject programs and detected 47\% of the existing equivalent mutants by applying this heuristic.
Similarly, Holling et al. \cite{holling2016nequivack,papadakis2012mutation} presented an approach for identifying non-equivalent mutants and improving the confidence of the mutation score. By using static analysis and symbolic execution they defined a six-steps procedure to determine which mutants are \textit{non-equivalent}. \textit{Non-equivalent} mutants are identified every time they find a counter-example input for which the outputs of a pair of functions (the original function and the mutant one) is different. In the case no counter-example is found, then the mutant is classified as \textit{unknown}. 
In like manner, Riener et al. \cite{riener2011test} proposed the Symbolic Bounded Model Checking (SymBMC) procedure for the automated generation of test cases from a set of mutants. The approach examines the original program and its mutants on the same input and seeks for executions resulting in different observable output for the program and mutants, if the procedure finds a failing execution, the input data is saved as an effective new test case. Every time a new test case is found, the mutant is defined as non-equivalent, in a similar way to Holling's approach \cite{holling2016nequivack}.

Other approaches consider the use of software clones (i.e., similar code fragments) to detect equivalent mutants \cite{kintis2013identifying}, in this work Kintis proposed that mirrored mutants (i.e., mutants that belong to the same software clone) present the same behavior with respect to each other. So, for a set of mirrored mutants, is enough to prove equivalence of one mutant, instead of trying to detect equivalence for the whole set.

With a different approach, Adamopoulos et al. \cite{adamopoulos2004overcome} introduced a co-evolutionary technique for detecting equivalent mutants. The technique defines a fitness function that sets a poor fitness value to an equivalent mutant. Through the fitness function equivalent mutants are removed during the co-evolutionary process, and only mutants that are hard to kill and test cases that are good at detecting mutants are kept for future iterations of the algorithm. On the other hand, Maldonado et al. \cite{maldonado2005bayesian} developed a Bayesian Learning-Based technique for helping tester to detect equivalent mutants using an inference algorithm.

% 	\item Margrave's change-impact analysis \cite{martin2007fault} (READ)
% 	\item Using Lesar model-checker for eliminating equivalent mutants \cite{du2008towards} (READ)

% \textbf{Avoiding equivalent mutant generation techniques}

% \begin{itemize}
% 	\item \textbf{Selective mutation} \cite{mresa1999efficiency}:
% 	Randomly selecting 10\%, 20\%, 30\%, 40\%, 50\% and 60\% of the mutants results in a fault loss of approximately 26\%, 16\%, 13\%, 10\%, 7\% and 6\% respectively \cite{papadakis2010empirical}.

\subsubsection{Reducing Equivalent Mutant Techniques}

The second group of techniques aims to reduce the amount equivalent mutants produced during mutation process.
In this direction, Gr\"{u}n et al. \cite{grun2009impact} proposed that mutants that do not alter the dynamic control-flow with respect to the original program, should be avoided since they have a greater chance of being equivalent. To measure the dynamic control-flow difference (i.e., the impact), they defined the impact as the number of classes with different statement coverage. On the empirical evaluation performed on JAXEN (12,449 LOC), by analyzing the mutants that were covered by the test bench but not killed, they performed several experiments with two main results. 
The first result was about the relation between impact on control-flow and likelihood of the mutant of being non-equivalent, the experimental results were obtained by randomly selecting 20 mutants, they discovered that the 60\% of the mutants with control-flow impact were classified as non-equivalent, and that the 60\% of the mutants without control-flow impact were classified as equivalent. 
The second result was about the relation between the impact on control-flow degree and likelihood of the mutant of being non-equivalent, the results were obtained by analyzing two specific subsets (e.g., (a) subset with the 20 mutants with highest impact, and (b) subset with the 20 mutants with lower impact), from the first subset with higher impact, 90\% of the mutants were non-equivalent, while on the second subset 55\% were equivalent mutants. The conclusion is that testers can effectively focus on mutations with higher impact, at the expense of loosing mutations that could reveal another type of faults.

In the same research line, Schuler et al. \cite{schuler2009efficient} demonstrated that mutants that violate dynamic invariants (i.e., invariants derived from data collected using dynamic analysis) are less likely to be equivalent and should be preferred over those that do not alter invariants with respect to the original program version. In an empirical evaluation performed on JAXEN, they analyzed two subset of mutants (e.g., (a) 12 random mutants that do not violate invariants, and (b) 12 mutants with highest score of invariants violated during execution), the results indicate that from the violating mutants 83\% were non-equivalent, while from the non-violating mutants only 33\% were non-equivalent.

A different perspective is to focus in code coverage of mutants, for example, Schuler et al. \cite{schuler2010covering,schuler2013covering} discovered that mutants that change code and data coverage with respect to the original program, has a likelihood of 68\%-79\% to be non-equivalent.
On an empirical evaluation performed on 140 mutations from seven open-source projects, the authors discovered (1) that operators that modify the control-flow such as \textit{negate jump condition} or \textit{omit method call} produce less equivalent mutants (30\%) than operators that only change data (57\%), and (2) that from several metrics based on differences in statement (i.e., number of methods that have at least one statement that is executed at a different frequency between mutated and original program) and data (i.e., number of methods that have at least one different return value between mutated and original program) coverage, the authors discovered that by measuring both impact on data and code coverage is possible to obtain few false negative results (61\%) when assessing if a mutant is equivalent or not, in comparison with measuring data and code coverage independently (56\% and 67\% ,respectively).

Program slicing has been presented as a possible mechanism for reducing equivalent mutants \cite{voas1997software, hierons1999using, harman2001relationship}: in particular, Harman et al. \cite{harman2001relationship} presented a technique for reducing equivalent mutation generation using program dependence analysis. The idea is to avoid mutants that fail to propagate corrupted data into the inspection set at the probe point, a mutant that fails to propagate specific data means that no semantic change is being introduced on the software behavior. To carry on this analysis, the authors used a method called \textit{JR-dependence}, the method allows to relate variable and node pairs rather than simply considering nodes. With \textit{JR-dependence} is possible to know the set of variables that can and cannot be used to kill a mutant, which is beneficial for mutation testing. 
Consequently, Offutt et al. \cite{offutt2006class} used the guidelines by Harman et al. \cite{harman2001relationship} to eliminate equivalent mutants by applying the optimizations directly in the implementation of object oriented operators developed on the MuJava mutation testing tool. 

On the other hand, Kintis et al. \cite{kintis2014using,kintis2015medic} discovered that through static data-flow analysis is possible to reveal code locations that are sensitive to produce equivalent mutants together with the mutation operators being applied. The first pattern is called \textit{Use-Def} Problematic Pattern and seeks for uses of variables that reach a definition and can be mutated by a mutation operator that produces in-place changes (e.g., $\{x = (m + x)/2\}$ mutate to $\{x = (m + x++)/2\}$), specifically this pattern can be applied on definitions at the same line, basic block and between different basic blocks. The second pattern is called \textit{Def-Def} Problematic Pattern and seeks for subsequent definitions of a certain variable, in other words, if a definition reach another definition of the same variable without a prior use of the variable, then any mutation to the first definition cannot be revealed since its being redefined. On the experimental evaluation guided on 6 real-world programs, the technique was able to automatically detect 58 out of 84 equivalent mutants, leaving only a 30\% of mutants to be analyzed manually.

Research has proved that higher order mutation testing can be helpful in reducing equivalent mutants \cite{jia2009higher,kintis2010evaluating,offutt1992investigations,papadakis2010empirical}, since two or more mutations are applied simultaneously, the chances of producing equivalent mutants decreases consistently. For instance, Papadakis and Malevris \cite{papadakis2010empirical}, worked on a approach for higher order mutants for the C programming language that lead to a reduction of approximately 80-90\% of the generated equivalent mutants, with a fault detection ability loss only of 11-15\%. 
% For instance, Offutt demonstrated that the set of test data developed for first order mutants (FOMs) actually killed a higher percentage of mutants when applied to second order mutants (SOMs) \cite{offutt1992investigations}. 
% Jia and Harman identified six different types of HOMs \cite{jia2009higher} and presented a categorization of HOMs. They introduced the concept of subsuming and strongly subsuming HOMs.
% Polo et al. \cite{polo2009decreasing} studied three strategies to combine FOMs and generate mutants, and found that they can achieve significant cost reductions without losing any effectiveness (they reduced the number of mutants in a approximately 50\%, without much decrease in the quality of the test suite).
%Instead, Kintis et al. \cite{kintis2010evaluating} developed a solution for the Java language, they state that SOMs achieve higher collateral coverage for strong mutation as compared with third or higher order mutants. With their approach they obtained a mutant reduction of between 65-87\% and a loss of test effectiveness from 1.75-4.2\%.
% Mateo et al. \cite{mateo2012validating,madeyski2013overcoming} found that second order mutants (SOM) are significantly more efficient that first order mutants (FOM).

\endinput

\begin{table*}[ht]
\centering
\scriptsize
\begin{tabular}{lllllllp{4cm}}
\toprule
Author(s)          & Year   & Language & \begin{tabular}[c]{@{}l@{}}Largest\\Subject\end{tabular} & \begin{tabular}[c]{@{}l@{}}\#Eq. \\ Mutants\end{tabular} & \begin{tabular}[c]{@{}l@{}}Available \\ Tool\end{tabular} & Category                                                 & Findings                                                                                      \\
\midrule
Baldwin \& Sayward \cite{baldwin1979heuristics} & 1979   &          &                                                           &                                                          &                                                           & Detect                                                   & Compiler optimization can be used to detect equivalent mutants                                \\
Acree  \cite{acree1980mutation}       & 1980   & Fortran  &                                                           & 25                                                       &                                                           & Detect                                                   & Human make mistakes when they identify equivalent mutants                                     \\
Offutt \& Craft \cite{offutt1994using}   & 1994   & Fortran  & 52                                                        & 255                                                      &                                                           & Detect                                                   & Compiler optimisation can detect on average 45\% of equivalent mutants                        \\
Offutt \& Pan \cite{offutt1996detecting,offutt1997automatically}     & 1996-7 & Fortran  & 29                                                        & 695                                                      & Yes                                                       & Detect                                                   & Constraint-based testing can detect on average 47\% of equivalent mutants                     \\
Voas \& McGraw \cite{voas1997software}    & 1997   &          &                                                           &                                                          &                                                           & Detect                                                   & Slicing may be helpful in detecting equivalent mutants                                        \\
Hierons et al. \cite{hierons1999using}      & 1999   &          &                                                           &                                                          &                                                           & \begin{tabular}[c]{@{}l@{}}Detect/\\ Reduce\end{tabular} & Program slicing can be used to detect and assist the identification of equivalent mutants     \\
Harman et al.  \cite{harman2001relationship}    & 2001   &          &                                                           &                                                          &                                                           & \begin{tabular}[c]{@{}l@{}}Detect/\\ Reduce\end{tabular} & Dependence analysis can be used to detect and assist the identification of equivalent mutants \\
Adamopoulos et al  & 2004 \cite{adamopoulos2004overcome}  &          &                                                           &                                                          &                                                           & Reduce                                                   & Co-evolution can help in reducing the effects of equivalent mutants                           \\
Grun et al. \cite{grun2009impact}       & 2009   & Java     & 12,449                                                     & 8                                                        & Yes                                                       & Reduce                                                   & Coverage Impact can be used to classify killable mutants                                      \\
Schuler et al. \cite{schuler2009efficient}    & 2009   & Java     & 94,902                                                     & 10                                                       & Yes                                                       & Reduce                                                   & Invariants violations can be used to classify killable mutants                                \\
Schuler \& Zeller \cite{schuler2010covering,schuler2013covering} & 2010-2 & Java     & 94,902                                                     & 63                                                       & Yes                                                       & Reduce                                                   & Coverage impact can be used to classify killable mutants                                      \\
Nica \& Wotawa \cite{nica2012using}    & 2012   & Java     & 380                                                       & 1,424                                                     &                                                           & Detect                                                   & Constraint-based testing can detect equivalent mutants                                        \\
Kintis et al. \cite{kintis2012isolating,kintis2015employing}     & 2012-4 & Java     & 94,902                                                     & 89                                                       &                                                           & Reduce                                                   & Higher order mutants can be used to classify killable mutants                                 \\
Kintis \& Malevris \cite{kintis2014using} & 2014   & Java     & 25,909                                                     & 84                                                       &                                                           & Detect                                                   & Data-flow patterns can detect 69\% of the equivalent mutants introduced by the AOIS operator  \\
Papadakis et al. \cite{papadakis2014mitigating}    & 2014   & C        & 513                                                       & 5,589                                                     &                                                           & Reduce                                                   & Coverage impact can be used to classify killable mutants                                      \\
Papadakis et al. \cite{papadakis2015trivial}    & 2014   & C        & 362,769                                                    & 9,551                                                     & Yes                                                       & Detect                                                   & Compilers can be used to effectively automate the mutant equivalence detection               \\
\bottomrule
\end{tabular}
\end{table*}



% !TEX root = MutationTestingSurvey.tex

\subsection{Solutions to Minimize Redundant Mutants}
\label{sec:opt:redundant}

The term \INDEX{redundant mutant} is used to refer to mutants that show the same behaviour of other mutants, i.e., they
are killed every time other mutants are also being killed. 
%Redundant mutants are another particular category of mutants, these mutants do not contribute to the testing process, since they are killed every time other mutants are also being killed. 
\emph{The main drawback of redundant mutants is that they can artificially inflate the apparent ability of a test technique to detect faults, in other words they tend to skew the mutation score measurement leading to serious threats to the validity of empirical research}~\cite{papadakis2016threats}.

%\DONE{Is it possible to add an example?}

% !TEX root =  ../MutationTestingSurvey.tex

\noindent\begin{minipage}{.45\textwidth}
\begin{lstlisting}[style=CStyle, caption=Excerpt from mutant M4 of isPalindrome., label=redudantexample1]
while (h > l) 
{ 
	...
	//l++;
	h--;
}
\end{lstlisting}
\end{minipage}\hfill
\begin{minipage}{.45\textwidth}
\begin{lstlisting}[style=CStyle, caption=Excerpt from mutant M5 of isPalindrome., label=redudantexample2]
while (h > l) 
{ 
	...
	l++;
	// h--;
}
\end{lstlisting}
\end{minipage}

In Section~\ref{sec:process} we introduced a mutation testing example for the function \texttt{isPalindrome}. 
Listing~\ref{redudantexample1} and~\ref{redudantexample2} show excerpts from mutants \textit{M4} and \textit{M5} obtained with the \textit{SSDL} operator. In this case, both mutants are considered not equivalent with respect to the original program, but they are redundant between each other, because they are being killed by the same test cases, that is, the test cases exercising the inputs \texttt{abba} and \texttt{aba}.
\REVTWO{C32}{In this case, the solution to the problem is dual. The first solution accounts for introducing one new test case that should produce different outputs for both mutants, which in this case is infeasible since there is no test case able to produce different output for mutants in Listing~\ref{redudantexample1} and~\ref{redudantexample2}. The second solution consists of excluding randomly one of the two mutants.}

%\DONE{The following sentence is not good. Can we say something more, for example how many SUT they considered, which mutation operators they considered?}

To highlight that redundant mutants are a recurrent problem, Kintis et al.~\cite{kintis2010evaluating} showed in an experiment that 9\% of mutants were redundant. For the experimental evaluation the authors considered 15 SUT, on a total of 372 LOC, and 6\,127 test cases, and applied the sufficient set of operators for generating the mutants. 

\MREVISION{C15}{Identifying all the redundant mutants is equal to finding the minimal set of mutants. \textit{A minimal set of mutants has no redundancy, that is, every test set that kills these mutants will also kill all remaining mutants}~\cite{kurtz2015static}.} 
%Similar to the case of equivalent mutants, finding the minimal set of mutants is an undecidable problem, and it can be only approximated.}

%\DONE{The definition below is not clear. I cannot understand teh difference between them}
In the literature, redundant mutants are divided into two categories. The first is the category of \INDEX{duplicated mutants}, that is, mutants that are equivalent with each other but not equivalent to the original program. The second category concerns \INDEX{subsumed mutants}, that is, mutants that are not equivalent with each other but are killed by the same test cases. 
\MREVISION{C15}{Only duplicate mutants are a problem for test suite evaluation; indeed, subsumed mutants are mutants that capture different faults. Unfortunately, automatically distinguishing the two cases is not feasible since it reduces to the problem of automatically identifying duplicated mutants, which, in turn, coincides with the problem of identifying equivalent mutants. Often, automated techniques simply identify mutants that lead to the same test failures, i.e., both duplicated and subsumed mutants.}

Previous studies showed redundancies between mutation operators and proved that a certain subset of operators is sufficient to measure test effectiveness. For instance, Rothermel et al.~\cite{rothermel1996experimental} and then Andrews et al.~\cite{andrews2005mutation} proposed a small set of operators that is lead to a sufficiently accurate approximation of the results obtained by using all possible operators (e.g., replace numerical constant, negate jump condition, replace arithmetic operator, omit method calls). In the same direction, Namin et al.~\cite{siami2008sufficient} proposed a statistical analysis procedure for identifying a subset of operators that could predict mutation score, their approach reduced mutants in a 93\% on C programs. 

More recently, Delamaro et al.~\cite{delamaro2014designing} designed deletion operators, and found that they form a cost-effective alternative to other operators, i.e., they produce less redundant mutants. The deletion operator by itself has been proven to be the most effective for fault detection~\cite{delamaro2014designing}.

%With a different perspective, 
Papadakis and Malevris~\cite{papadakis2012mutation} and then Kurtz et al.~\cite{kurtz2015static} proposed a path selection strategy (i.e., they generate new test inputs) for selecting the test cases able to effectively kill mutants using \INDEX{symbolic execution}, and to consequently decrease the number of redundant mutants. 
The authors suggest constrained versions of the logical, relational and unary operators for generating less redundant mutants. 
In a similar manner, Just et al.~\cite{just2012redundant,just2015higher} proved that these three operators are better at detecting faults that the rest of mutation operators.

Delgado et al.~\cite{delgado2017assessment} show that some operators naturally produce more redundant mutants than others and
developed a selective approach for reducing the number of mutants without loss of effectiveness for C++ programs. 

The approach introduces a degree of redundancy for every mutation operator that helps developers to choosing the mutation operators with a lower degree of redundancy, based on the test cases defined in the project.

Another solution to reduce the number of redundant mutants is the application of trivial \INDEX{compiler optimisations}~\cite{papadakis2015trivial, kintis2017detecting,papadakis2019mutation}. 
It can identify duplicate mutants by comparing the optimised object code of each mutant. The empirical study guided by Kintis et al.~\cite{kintis2017detecting} showed that by using compiler optimisations it is possible to reveal 21\% and 5.4\% of C and Java mutants, respectively.


Finally, Shin et al. suggest to avoid discarding redundant mutants but, instead, augment the test suite with additional test cases so that 
each mutant can trigger a test failure that cannot be observed with other mutants~\cite{Shin:TSE:DCriterion:2018}. 
They introduce the \INDEX{distinguishing mutation adequacy criterion} to characterize test suites in which every mutant triggers a test failure that is not observed with other mutants.
Empirical results show that test suites that satisfy the distinguishing mutation adequacy criterion have a higher
 fault detection effectiveness than test suites that simply satisfy mutation coverage.





% !TEX root = MutationTestingSurvey.tex
\clearpage
\section{Building Blocks of the Code-driven Automated Test Suites Augmentation Process}
\label{sec:testGeneration}

This section describe the approaches that can be adopted to automatically generate test cases that kill mutants.
%Fabrizio:
To kill a mutant we need test cases that reach the mutation point (i.e., execute the mutated code), cause 
corruptions
%changes, or corruptions, 
in the program state right after the mutated code in comparison to the non-mutated version of the program, 
and manifest these corruptions into the program output 
%Fabrizio: an assertion is not an output
(e.g., by producing an erroneous value in a state variable verified by a test assertion) 
thus leading to a software failure~\cite{papadakis2019mutation}, these conditions are also known as the \INDEX{killing conditions} of a mutant.

In the literature, there exist two groups of approaches for 
%driving automated test generation processes aimed to kill mutants, based on the above mentioned conditions. 
generating test cases that kill mutants:
%Fabrizio: "symbolic procedures" is really the right terminology?
%Oscar: another term might be "Constraint-based"
%anywa, below symbolic execution is just part of the solutions adopted so we need a broader term
%The first group consists of techniques based on symbolic execution; the second group relies on search-based test generation techniques.
approaches based on constraint-programming, and approaches based on evolutionary computation.

%\subsection{Symbolic Test Generation Procedures}
\subsection{Test Generation based on Constraint Programming}

Techniques based on constraint programming use some form of automated reasoning (e.g., Propositional Satisfiability or Constraint Solving~\cite{SATandCPsurvey:2006}) to derive data that satisfy all the conditions necessary to kill a mutant~\cite{offutt1997automatically}.

One possibility consists of generating constraints that capture the three killing conditions of a mutant: \INDEX{reachability} (i.e, the test case should execute the mutated statement), \INDEX{necessity} (i.e., the test case should cause an incorrect intermediate state if it reaches the mutated statement), and \INDEX{sufficiency} (i.e., the final state of the mutated program should differ from that of the original program)~\cite{offutt1997automatically}. Existing approaches, differ for the strategy adopted to automatically generate these constraints from the program under test.

Offutt et al.~\cite{offutt1997automatically}, for example automatically derive such constraints from the program by extracting the predicate expressions on the program's control flow graph.
Then, such constraints are encoded to form a constraint system, in their approach, they propose three strategies for identifying infeasible constraint systems, the contradictions to such systems are the new test cases for the program under analysis.

% !TEX root =  ../MutationTestingSurvey.tex

\begin{lstlisting}[style=CStyle, caption=midval function returns the mid value between three integers., label=midval, mathescape=true]
int midval (int x, int y, int z) {
	int midval;

	midval = z;
	if (y < z) {
		if (x < y) {
			midval = y;
		}
		else if (x < z) 
 $\Delta$ else if (x <= z) {
			midval = x;	
		}
	}
	else 
		if (x > y) {
			midval = y;
		} else if (x > z) {
			midval = x;
		}
	return midval;
}
\end{lstlisting}

In the following, we introduce an example of the application of Offutt's~\cite{offutt1997automatically} approach by using the \texttt{midval} function presented in Listing~\ref{midval}. This function has a mutation on line 9, which has been mutated into line 10. According to their approach, the three killing conditions would be the following:

\begin{itemize}
	\item Reachability $C_R: (y < z) \wedge (x \geq y)$
	\item Necessity $C_N: (x < z) \neq (x \leq z)$
	\item Sufficiency $C_S:$ Output(P) $\neq$ Output(M)
\end{itemize}

$C_R$ defines the condition required to reach the mutated statement, in this case the conjunction between the predicate of the first if condition and the negation of the second if condition. $C_N$ defines the condition required to assure a different program state between the original and mutated version of the program right after the mutation point. Finally, $C_S$ defines the condition necessary to demonstrate that both the original and the mutated program returns different values.

Holling et al.~\cite{holling2016nequivack} proposed to use a \INDEX{symbolic execution} approach to identify new test cases, the idea is to execute symbolically both the original and the mutated function, and then to check if their return values are equivalent or not. 

% Symbolic execution

Symbolic execution determines what inputs cause each part of a function to execute. To symbolically execute a function, it is necessary to replace the original inputs (i.e., concrete values) with symbolic ones. The \INDEX{symbolic values} represent a set of possible concrete values that lead to a certain program path (i.e., path condition). 

In the following, we introduce an example of symbolic execution applied to the function \texttt{f} presented in Listing~\ref{symbolic}, the function \texttt{f} reads in a value and fails if the input is 6.

\begin{lstlisting}[style=CStyle, caption=Symbolic Execution Example., label=symbolic]
int f() {
	int y, z;
	y = read();
	z = y * 2;
	if (z == 12) {
		fail();
	} else {
		printf("OK");
	}
}
\end{lstlisting}


During a concrete execution, the function \texttt{f} would read a concrete input value (e.g., 5) and assign it to \texttt{y}. The execution would then proceed with the multiplication and the conditional if branch, which would evaluate to false and print \emph{OK}.

Instead, during symbolic execution, the function \texttt{f} reads a symbolic value (e.g., $\omega$) and assigns it to \texttt{y}. The function would then proceed with the multiplication and assign $\omega * 2$ to \texttt{z}. When reaching the if statement, the program would evaluate the expression $(\omega * 2 == 12)$. 
At this point of the execution, $\omega$ could take any value, and symbolic execution can therefore proceed along both branches, by creating two paths. 
Each path is associated to a path constraint and a copy of the program state.
%Fabrizio: the following is misleading in the example because we do not have it...
% a copy of the program state at the branch instruction and . 
In this example, the path constraint is $(\omega * 2 == 12)$ for the if branch and $(\omega * 2 != 12)$ for the else branch. 
Both paths can be symbolically executed independently. When a path terminates (e.g., as a result of executing the function \emph{fail()} or simply when it reaches the return instruction), the symbolic execution process computes a concrete value for $\omega$ by solving the path constraints accumulated on each path. 
These concrete values can be used as inputs for concrete test cases that cover the paths. In this example, the constraint solver would determine that in order to reach the fail() statement, $\omega$ needs to be equal to 6.
%%

Holling's approach~\cite{holling2016nequivack} to automatically identify equivalent mutants relies on the observation that 
%Fabrizio: the followng is very convoluted, it says the same thing in two different ways
%when both the mutated and the original function return always the same outputs then 
%means that there are not concrete values making the mutated function to produce a different output.
two mutants are equivalent when there are no concrete values making the mutated function produce an output that is different from the one of the original function.
If a value that makes the two functions generate distinct results can be found, the mutant is non-equivalent.
%Fabrizio: also the following is not understandable
%Each non-equivalent value proves that a new input that makes the original and the mutated function generate different results, has been found. 
To automate the generation of inputs, Holling et al. rely on KLEE~\cite{cadar2008klee}, a symbolic execution tool for programs written in C/C++.

We introduce an example of Holling's approach in Listing~\ref{function}. The top part of Listing~\ref{function} shows the function \texttt{isPositive}, which checks if an integer number is positive or not. The bottom part of Listing~\ref{function} presents the mutated version of \texttt{isPositive}, where the relational operator $\geq$ has been replaced by the operator $>$.
To automate the generation of inputs using KLEE, all the parameters need to be treated as symbolic values.
%F: "the \texttt{isPositive} parameter \texttt{num}" it's impossible to understand what is the parameter and who owns the parameter
%In Holling's approach, first, the \texttt{isPositive} parameter \texttt{num} needs to be treated as symbolic value, which is done by the function \texttt{make\_symbolic} (see Listing~\ref{example}) which converts concrete variables to symbolic by considering its memory address and size. 
This is achieved by function \texttt{make\_symbolic} (see Listing~\ref{example}) which converts concrete variables to symbolic ones by considering their memory address and size. 
%F: please fix
In Listing~\ref{example}, the parameter \texttt{numSymbolic} is made symbolic in Line 3. Then, the original and mutated functions are called using the symbolic arguments in Lines 5 and 6.
%Finally, the \texttt{assert} function of line 8 verifies if the integer return values of both functions are equal or not.
Finally, we need to introduce an assertion that makes the symbolic execution engine \EMPH{look for inputs that make the output of the two functions different}. 
In Listing~\ref{example}, this is achieved with an assertion that verifies that the return values of the two functions the same (see Line 8). 

Despite being counterintuitive, this approach is effective because symbolic execution engines aim to identify inputs that falsify the assertions in the program. 
When the equality is falsified, then the two functions can produce a different output for a same input\footnote{In the following sections, we will show that this same criterion also indicates that the two functions are non-equivalent}.
The input that falsifies the equality can thus be used to improve the test suite enabling it to kill the mutant.
% make the original an mutate function geenrate distinct results
% of identifying inputs that lead to different output falsifying the assertion.

%In the concrete, if we consider the implementation of \texttt{isPositive} and \texttt{MUT\_isPositive}, 
In the example of Listing~\ref{example}, KLEE will indicate that the return values of the original and mutated function differ when \texttt{num} is equal to zero.
A new test case exercising function \texttt{isPositive} with \texttt{num=0} should thus be added to the test suite in order to kill the mutant.

\begin{lstlisting}[style=CStyle, caption=isPositive and MUT\_isPositive functions, label=function]
int isPositive(int num){
	if (num >= 0){
		return 1;
	} else {
		return 0;
	}
}

int MUT_isPositive(int num){
	if (num > 0){
		return 1;
	} else {
		return 0;
	}
}

\end{lstlisting}

\begin{lstlisting}[style=CStyle, caption=Holling's approach for test case generation., label=example]
void test () { 
	int numSymbolic; 
	make_symbolic(&numSymbolic, sizeof(numSymbolic), "numSymbolic"); 
 	
 	int original_ret = isPositive(numSymbolic); 
	int transformed_ret = MUT_isPositive(numSymbolic); 
 
	assert(original_ret == transformed_ret); 
}
\end{lstlisting}

Similarly to Holling's approach, Riener et al.~\cite{riener2011test} proposed to use \INDEX{bounded model checking} techniques to search for these counter examples.
In their bounded model checking approach, the original program and the mutant are unrolled with respect to a certain maximum bound. In program unrolling, loops are re-written as a repeated sequence of similar independent statements. Then, both unrolled programs are encoded into a logic formula over the same input variables. To ensure that the mutation affects the output of the mutant, a propagation condition is encoded and added to the previous logic formula, the condition asserts that there exist at least one pair of different outputs under the assumption of equal inputs. In the last step, the formula is processed by a SMT-solver, if the solver finds a satisfying assignment, the inputs of the formula are translated into a new test case for the current program under analysis.

%F: you should explain what is bounded-model checking. If Riener do not provide background you can check my ISSTA'14 paper (Verification-Aided Regression Testing), there is a background section (2) on bounded model checking that starts with "The idea in BMC is to represent ..."

Compared to symbolic execution, one advantage of bounded model checking is that it does not require to execute the whole program but may focus on the mutated functions only.

To reduce the time required by the symbolic execution process, which needs to be performed against all the mutants of the software, Papadakis et al.~\cite{papadakis2011automatically, papadakis2010towards} propose to combine symbolic execution techniques and \INDEX{mutant schemata} to automatically generate test cases targeting the killing conditions induced by the different mutants embedded into the same executable. The approach targets \INDEX{weak mutation} testing and may not generalize to strong mutation testing. Indeed, ensuring the sufficiency property (i.e., verify that changes are propagated to outputs) for multiple mutants might lead to scalability issues not addressed by the proposed approach.

\subsection{Test Generation based on Evolutionary Computation}

Test generation approaches based on \INDEX{evolutionary computation} typically rely on population-based meta-heuristic optimization algorithms~\cite{harman2011strong}. 
They search for program inputs that could kill mutants under the guidance of a fitness function~\cite{harman2011strong}. 
The main research contribution of these methods is the definition of fitness functions that captures the killing conditions of a mutant and identifies test inputs that satisfy those conditions.

The \INDEX{fitness function} captures the killing conditions of a mutant. For instance, Ayari et al.~\cite{ayari2007automatic} proposed to use an evolutionary approach based on ant colony optimization (ACO) for automatic test input data generation on mutation testing. The ACO is an optimization algorithm inspired by the behavior of real ants, it is based on the ants ability to find the shortest path between their nest and the food source. In the study by Ayari et al.~\cite{ayari2007automatic}, the approach takes an existing test case and produces a new test case by slightly modifying its inputs. 
%F: you have to provide more details what does it mean "close"?
The fitness function measures the distance between the mutated statement, and the statement reached by the new test case (e.g., the reachability condition). More precisely, the distance is defined as the number of basic blocks between the two statements in the program's control flow graph.
Papadakis et al.~\cite{papadakis2011automatically}, instead, rely on fitness functions that capture the distance between mutated statement and the statement covering the branches of the different mutations (e.g., the necessity condition).

Fraser and Arcuri~\cite{fraser2015achieving} propose to use distinct distance metrics tailored to the specific operator used to generate the mutants.
%F: we need more examples for multiple operators. Also, you should clarify why distinct metrics are needed
This tailoring is needed because the necessity killing condition relies on changes in the program state and the execution of a mutated statement does not guarantee that the program state had been changed (i.e., values on the stack are different at the mutation point).
%Fabrizio: it was not comprehensible
%mutation operators change the program state, and in other cases the program state remains unchanged. Because of this, distinct metrics for measuring distance for each operators needs to be defined.
%\DONE{I cannot undertsand the following, we can ignore it.}
For example, the \textit{deletion operator}, which removes a statement, can or cannot change the program state. In case the mutation effectively changes the program state the distance is set to 0, otherwise the given value is 1.
% For example if the \textit{deletion operator} changes the program state (i.e., values on the stack are different at the mutation point) the distance is 0, otherwise the given value is 1. 
In the case of the \textit{insert unary operator}, which adds or subtracts 1 to a numerical value, the operator always change the program state, so the distance is set to 0 when the statement is reached. 
% For this reason, the mutants produced by this operator always affect the program state, so the distance is always 0 when the statement is reached.

Instead, in the case of the \textit{replace variable operator}, which replaces a specific variable with all other variables of the same type in the program scope, the distance is set to 0 only if the values of the variables being exchanged are different before executing the statement, otherwise it is 1.

\subsubsection{\REVTWO{C37}{Generation of test oracles}}
\label{sec:oraclesGeneration:codeDriven}

The automated generation of test oracles is a research topic that goes beyond the specific needs of mutation testing~\cite{Barr:Oracles:15,OLIVEIRA:Oracles:2014}.

Fraser et al. provide an overview of existing approaches for the automated generation of test oracles that have been integrated into existing test case generation tools~\cite{fraser2011mutation}.
A common solution consists of the automated synthesis of assertions for the test case, i.e., the generation of assertions  that reflect the output generated by the function under test when it is exercised with the input automatically generated by the software testing automation toolset. 
For example, Line~6 in Listing~\ref{isPositiveOracle} shows the oracle that can be automatically generated for the function analyzed in Listing~\ref{example} (i.e., function \emph{isPositive}). The oracle, in this case, consists of an assertion verifying that the value generated by function \emph{isPositive} matches the value '1', which is the value observed during test generation for the input value '0'. This is the approach implemented by Riener et al.~\cite{riener2011test}, who generate assertions that verify variables that present different values between the original and the mutated executions.

\begin{minipage}{15cm}
\begin{lstlisting}[style=CStyle, caption=Automatically generated oracle for function 'isPositive'., label=isPositiveOracle]
void test () { 
	int numSymbolic = 0;
 	
 	int ret = isPositive(numSymbolic); 
	
	assert( ret == 1 ); 
}
\end{lstlisting}
\end{minipage}

Randoop~\cite{PachecoLEB2007} allows annotation of the source code to identify observer methods to be used for assertion generation. Orstra~\cite{Xie:2006} generates assertions based on observed return values and object states.
DiffGen~\cite{Taneja:2008} extends the Orstra approach to generate assertions from runs on two different program versions.

A well known limitation of automatically generated assertions that reflect the actual values observed during execution is that
they need to be validated. More precisely, we need to ensure that the values expected by the assertions do not reflect a failure triggered by the test case (e.g, an erroneous value being returned). Such validation activity is typically performed manually by the engineers because it should be based on domain knowledge and system specifications. 
Specifications are generally written in natural language because, to reduce development costs, only few components of the system are specified using formal languages. For this reasons the automated verification of such assertions is infeasible.

Approaches that support engineers in the analysis of generated oracles exists and might be considered to speed up the process~\cite{Staats2012,PastoreICSE2015}. For example, 
Staats et al.~\cite{Staats2012} identify the variables to be verified by oracles to maximize the fault finding potential of the testing process is maximized with respect to the cost of manually verifying the correctness of each generated assertions. First, they generate a collection of mutants from the SUT. Second, the test suite (automatically generated) is run against the mutants using the original system as the oracle. Third, they select the variables to verify in test oracles by focussing on those variables that show different values in the original and the mutated version.

ZoomIn.~\cite{PastoreICSE2015} automatically identifies suspicious assertions. These are assertions that verify data that is generated by functions showing anomalies during test cases executions. Anomalies are detected by automatically deriving pre- and post-conditions of the functions of the SUT based on the data recorded during the execution of a manually implemented test suite. Anomalies consist of function executions that violate such pre- and post-conditions. 
To rank unsafe assertions, ZoomIn takes into consideration the number of violations of pre- and post-conditions produced by the related program variables. However, not all the constraint violations are equally important. Since erroneous behaviors are not frequent in adequately tested software, ZoomIn weights each constraint according to the number of times the constraint has been violated. The frequently violated constraints are likely imprecise constraints that erroneously detect legal values as anomalous values, while constraints that are seldom violated are likely to carry useful information. ZoomIn captures this aspect with the uniqueness score. 
Also, ZoomIn associates the assertions with a score, the suspiciousness score, that represents the likelihood the assertion is wrong. The suspiciousness score of an assertion depends on both the number and the uniqueness scores of the related constraints that are violated. Intuitively the highly scored assertions, i.e., the most suspicious assertions, are the ones associated with several constraint violations with high uniqueness scores.
To keep the inspection effort low and the effectiveness high, developers are assumed to inspect only the top assertions in the ranking. Results show that inspecting the top five unsafe assertions is enough to discover several faults without wasting time inspecting too many correct assertions.

% measure how often each variable in the system reveals a fault in a mutant and?based on this information?we rank variable effectiveness in terms of fault finding. Finally, we estimate? based on this ranking?which variables to include in the oracle data for an expected value oracle. The underlying hypothesis is that, as with mutation-based test data selection, oracle data that is likely to reveal faults in the mutants will also be likely to reveal faults in the actual system under test. This oracle data selection process is completely automated and requires no manual intervention. Once this oracle data is selected, the tester defines expected values for each element of the oracle data. Testing then commences with a?hopefully?small and highly effective oracle.


%According to recent survey on the topic, test oracles can be divided in three categories:
%oracles defined though a \INDEX{specification language} (i.e.,  is a notation for defining a specified test oracle, which judges whether the behaviour of a system conforms to a formal specification).

%  test oracles can be derived (Section 5);
%  test oracles can be built from implicit information
%(Section 6); and

%In the context of code-driven metamorphic testing
%
%
%
%Staats2012

\endinput

\subsubsection{Symbolic execution test generation} % (fold)

\begin{itemize}

	\item Dynamic symbolic execution approaches: embed the mutant killing conditions within the executable program and guide test generation towards these conditions.

	\item Papadakis et al.~\cite{papadakis2011automatically, papadakis2010towards}:
	Embed mutant infection conditions (cause a corruption to the program state) within the schematic functions that are produced by mutant schemata. Basically, we have all mutants encoded in a single executable with their killing conditions. Then, using dynamic symbolic execution we can directly produce test cases that targets those infection conditions.
\end{itemize}




% !TEX root = MutationTestingSurvey.tex

\subsection{Run-time Scalability}
\label{sec:opt:execution}

\endinput



If we have $n$ mutants and a test suite of $m$ tests, we have to perform $n \times m$ program executions at maximum.
The main issue remains the scalability of the amount of runnings needed to cover all mutants.

\begin{itemize}
	\item execute only mutants that are reachable by test cases
	\todoinline{easy to reimplement}
	\item is possible to record with one execution all the mutants that can be infected by a test
	\item instead of executing every mutant with every test, it is possible to execute all mutants at once, by monitoring the coverage of the infection conditions
	\item When mutation testing targets strong mutation and the case of test suites take a long execution time, we can reduce the execution time by checking if every mutant achieve weak mutation.
	\todoinline{How easy it is to determine if we have weak mutation in C programs? I think that one of the solutions is linked to split stream.} 
	If weak mutation is not achieved, i.e., if the program state is not infected, there is no point to execute the test case till the end. We can simply report it as not killed and save time.
	\item criticality: since determining whether a test execution can terminate or not is an undecidable problem, heuristic solutions are needed. Such as establishing time thresholds for running executions.
	\todoinline{but if we face this problem, isn't it enough to check if a test execution is taking twice the expected time? I mean, if the test case does not fail in 3 times the original execution time we can report the mutant as not killed. But the solution to this problem might be to add a timeout inside the test suite.}
	\item take advantage of common executions parts (between the original program and mutants). Called as Split-Stream execution
\end{itemize} 

Strong, weak and firm mutation
\begin{itemize}
	\item Strong mutation:
	\item Weak mutation: the mutant is killed if any mutated component changes its state (variable reference, variable assignment, arithmetic expression, relational expression, boolean expression). Advantages: weak mutation does not need a complete execution. Cons: sacrifices test effectiveness with test effort.
	\item Firm mutation: the killing process occurs between weak and strong mutation
\end{itemize} 

Runtime optimization techniques
\begin{itemize}
	\item Interpreter based technique: results are analysed from the source code
	\item Compiler based technique: results are analysed from the execution of binary code 
	\todoinline{how much gain can it give?}
	\item Mutant Schema Generation: creates a metaprogram, there is research on mutating bytecode
	\todoinline{Does MUSIC implement it? We can implement it as a postprocessing. We compare files generated by MUSIC. We identify the changed lines (mutants). Then we generate a single file where we have options for each changed line.}
\end{itemize}

% !TEX root = MutationTestingSurvey.tex

\section{Reduce/Prioritise Test Cases}

\label{sec:reduceprioritise}

\begin{itemize}
	\item Reduce: removal of test cases that are somehow redundant.
		\todoinline{which type of redundancy the focus on?}
	\item Prioritise: ordering test cases in a way that mutants are killed as early as possible.
\end{itemize}





% !TEX root = MutationTestingSurvey.tex

\section{Reduce/Prioritise Test Cases}

\label{sec:reduceprioritise}

\begin{itemize}
	\item Reduce: removal of test cases that are somehow redundant.
		\todoinline{which type of redundancy the focus on?}
	\item Prioritise: ordering test cases in a way that mutants are killed as early as possible.
\end{itemize}



% !TEX root = MutationTestingSurvey.tex

\subsection{Mutation Score Calculation}
\label{sub:mutationscore}

%\DONE{Is there anything we can add here?}

\MREVISION{C6}{The \INDEX{mutation score} captures, in percentage points, the quality of a test suite. It measures the percentage of mutants that had been killed by the test suites.} 
The mutation testing process is driven by the mutation score; the process iterates multiple times until the mutation score reaches a certain threshold. 

According to recent studies, there is a relation between the mutation score and the fault revelation ability of mutation testing.
\MREVISION{C7}{A recent empirical evaluation, for example, has shown that \textit{achieving high mutation scores improves significantly the fault detection capability of a test suite}
~\cite{papadakis2018mutation}. 
\DONE{The following sentence is incomplete. You cannot say "similar fault detection ability that statement and branch coverage". Do you mean 100\% branch coverage (i.e., branch adequate)? Also putting together branch and statement is tricky, because branch is a stronger criterion. }
This evaluation shows that a test suite that reaches a mutation score of 80\% has a similar  similar fault detection capability of one that achieves 80\% branch coverage. 
\DONE{Why do we have to put both 90 and 95?}
Instead, a test suite that reaches a mutation score of 90\% outperforms this code coverage criterion. The drawback of this result, is, however, that the fault detection ability of mutation testing is high only with a high mutation score.}

%To reliably compute the mutation score it is necessary to identify equivalent and redundant mutants. 
As seen in Section~\ref{sec:opt:equivalent} and~\ref{sec:opt:redundant}, equivalent and redundant mutants can affect (i.e., cause an overestimation or underestimation) the mutation score~\cite{papadakis2016threats,kintis2017detecting}.
Hence, it is desirable to identify the presence of such mutants before estimating the fault revelation ability of a test suite.



\newpage

\bibliographystyle{IEEEtran}
\bibliography{bibliography}


\end{document}  